\documentclass{situationist-preprint}

\title{The Distribution of Animal Consciousness}
\msid{MND-23-a}
\msfamily{MND-23}
\msdiscipline{Philosophy of Mind}
\mstype{Principal manuscript}
\mskeywords{animal consciousness; sentience; markers; homology; Morgan's canon; birds; cephalopods; insects; distributed processing}

\begin{document}
\maketitle

\begin{abstract}
Which animals are conscious? The question cannot be settled by any theory of consciousness now available, since the theories disagree precisely where the animals differ from us. We defend a theory-light method: consciousness is ascribed to a group of animals when the best explanation of a convergent set of behavioral and neural markers, weighed by their evolutionary relation to our own case, is that its members have experiences. We argue that Morgan's canon, properly understood, does not tell against such ascriptions, because for animals closely related to us the hypothesis of shared consciousness is the more parsimonious one. Applying the method, we conclude that mammals and birds are conscious: the markers converge, and the relevant structures are homologous with our own. Cephalopods are probably conscious: the markers converge, but they arise in an independently evolved nervous system in which most processing is distributed through the arms, so the inference rests on analogy rather than homology. For insects the evidence remains inconclusive: some markers are present, others absent or untested, and simpler explanations of the positive findings have not been excluded. We draw the consequences for the ethics of animals.
\end{abstract}

\section{The problem}

Descartes held that animals are machines whose behavior, however complex, is produced without thought or feeling (\emph{Discourse on the Method} V). Few now accept the conclusion, but the problem it answered remains. We know that we are conscious; we infer that other human beings are, on grounds examined in MND-06; and we must decide, for creatures whose bodies and brains differ from ours by degrees that increase with evolutionary distance, whether the same inference extends to them. \citet{Nagel1974} took for granted that bats are conscious and asked what their experience is like; MND-03 answers that question. Here we ask the prior one.

The question is about phenomenal consciousness: whether there is something it is like to be the animal, whether it has experiences with qualitative character in the sense defended in MND-10. It is not about intelligence, which can come apart from consciousness, nor about self-consciousness, which many conscious animals may lack. It matters for ethics. ETH-07 holds that sentience, the capacity for valenced experience, is necessary and sufficient for moral status; ETH-10 that factory farming is impermissible. Both rely on an answer to the question addressed here.

Positions range from the restrictive to the expansive. \citet{Carruthers2019} argues that, on the best theory of consciousness, there is no fact of the matter whether animals other than humans are conscious, since the question has no determinate answer outside the human case. \citet{Key2016} argues that fish lack the cortical structures necessary for pain. At the other end, \citet{Barron2016} argue that insects have a midbrain structure that performs the integrative function that supports subjective experience in vertebrates, and are therefore conscious, and \citet{Feinberg2016} argue that consciousness arose early in vertebrate and arthropod evolution. Between these, \citet{Birch2020} and \citet{Birch2022} defend a marker-based method that we adopt with modifications, and \citet{Birch2024} develops it into a framework for policy under uncertainty. Public statements by researchers, such as the Cambridge Declaration on Consciousness of 2012 and the New York Declaration on Animal Consciousness of 2024, record a broad consensus for mammals and birds and a more guarded one for other groups.

\section{A short history}

The modern study of animal minds began with Darwin's argument for continuity between human and animal mental capacities and with the anecdotal method of Romanes. \citet{Morgan1894} reacted against the latter with his canon: in no case should an action be interpreted as the outcome of a higher psychical faculty if it can be interpreted as the outcome of one that stands lower in the psychological scale. The canon was taken by behaviorists as licence to eschew mental ascriptions altogether, and for half a century the scientific study of animal consciousness was largely abandoned. Its return came with cognitive ethology and, later, with neuroscience: the discovery that the subcortical structures implicated in affect are conserved across mammals, the argument of \citet{Merker2007} that consciousness does not require a cerebral cortex, and the reclassification of the avian brain, whose pallium is now understood as homologous with the mammalian pallium, including cortex, rather than as an enlarged basal ganglia \citep{Jarvis2005}. \citet{Edelman2005} proposed a list of neural and behavioral hallmarks of consciousness to be sought in non-mammalian species, and the marker approach has since developed from such lists.

\section{Framework}

\begin{definition}[Consciousness of a group]\label{d:cons}
A group of animals is \emph{conscious} if its typical adult members have phenomenally conscious experiences, at least some of which are valenced.
\end{definition}

The definition concerns typical adults; it leaves open the status of early developmental stages and of members with unusual neural conditions. The clause about valence restricts attention to the kind of consciousness that matters for ETH-07 and that the markers below are designed to detect.

\begin{definition}[Markers]\label{d:markers}
A \emph{marker} of consciousness is a behavioral or neural feature that, in humans, is associated with conscious rather than unconscious processing. We use the following families \citep{Birch2022,Gibbons2022}:
\begin{enumerate}
\item[(M1)] \emph{Nociceptive integration}: responses to noxious stimuli that are integrated in central structures rather than confined to reflexes.
\item[(M2)] \emph{Motivational trade-offs}: the animal weighs avoidance of a noxious stimulus against other needs, such as food, in a flexible way.
\item[(M3)] \emph{Wound-directed behavior}: tending, guarding or protecting an injured body part.
\item[(M4)] \emph{Analgesic and anxiolytic effects}: drugs that relieve pain or anxiety in humans alter the animal's behavior in the corresponding ways, including self-administration when injured.
\item[(M5)] \emph{Associative learning beyond reflex}: conditioned place avoidance, trace conditioning, and learning from novel noxious stimuli.
\item[(M6)] \emph{Affect-like states}: global states, such as optimism or pessimism biases, that modulate behavior across contexts.
\item[(M7)] \emph{Integrative neural architecture}: centralized or networked structures that integrate information from several sources and sustain it over time.
\end{enumerate}
\end{definition}

\begin{definition}[Homology and analogy]\label{d:hom}
A feature shared by two groups is \emph{homologous} if it is inherited from a common ancestor that had it, and \emph{analogous} if it evolved independently in each.
\end{definition}

\begin{principle}[Explanatory inference]\label{p:ibe}
One is justified in ascribing consciousness to a group to the degree that its being conscious is the best explanation of the markers its members show, compared with explanations that posit only unconscious processing.
\end{principle}

\begin{principle}[Convergence]\label{p:conv}
The more marker families a group shows, and the more independent their sources of evidence, the stronger the inference of Principle~\ref{p:ibe}; the absence of markers that have been sought and not found weakens it.
\end{principle}

\begin{principle}[Homology]\label{p:hom}
If the markers are produced by structures homologous with those that support consciousness in humans, the inference of Principle~\ref{p:ibe} is stronger than if they are produced by analogous structures, since in the homologous case the hypothesis that the structures do the same work in both groups is the more parsimonious one.
\end{principle}

Principle~\ref{p:hom} rests on an argument of \citet{Sober2005}. Morgan's canon tells us to prefer the lower faculty, but cladistic parsimony tells us to prefer the hypothesis that requires fewer evolutionary changes. If humans are conscious and a closely related group shows the same markers through homologous structures, the hypothesis that the group lacks consciousness requires that consciousness was gained in our lineage alone after the split, or lost in theirs; the hypothesis that it is shared requires no such change. For closely related groups the two parsimony principles conflict, and cladistic parsimony is the one with an evolutionary rationale. For distantly related groups with analogous structures, cladistic parsimony is silent, and the canon retains whatever force it has.

\begin{definition}[Evidential verdicts]\label{d:verdict}
A group is \emph{conscious} (for the purposes of judgment) if the inference of Principle~\ref{p:ibe} is strong by both Principle~\ref{p:conv} and Principle~\ref{p:hom}; \emph{probably conscious} if it is strong by Principle~\ref{p:conv} but rests on analogy; and the evidence is \emph{inconclusive} if the markers are partial or unconvergent and an explanation without consciousness has not been excluded.
\end{definition}

\section{Mammals}

\begin{lemma}\label{l:mammals}
Mammals are conscious.
\end{lemma}

\begin{proof}
Mammals show every marker family. Rodents and other mammals show integrated nociception (M1) and trade off avoidance of noxious stimuli against food and shelter (M2); they guard injured limbs (M3); injured rats self-administer analgesics, and analgesics and anxiolytics alter their behavior as they alter human behavior (M4); they learn conditioned place avoidance and trace associations (M5); they show judgment biases that track welfare (M6). The structures that support these capacities in humans, the thalamocortical system and the subcortical systems for affect, are homologous across mammals (M7). By Principles~\ref{p:conv} and \ref{p:hom}, the inference of Principle~\ref{p:ibe} is strong on both counts. By Definition~\ref{d:verdict}, mammals are conscious.
\end{proof}

\section{Birds}

\begin{lemma}\label{l:birds}
Birds are conscious.
\end{lemma}

\begin{proof}
Birds show integrated nociception, motivational trade-offs, injury guarding, and analgesic effects on pain behavior (M1--M4). They show sophisticated learning (M5), including in corvids the planning and caching behavior that has been compared with that of great apes \citep{Emery2004}, and judgment biases (M6). Their pallium is homologous with the mammalian pallium, though organized in nuclei rather than layers \citep{Jarvis2005}, and \citet{Nieder2020} found neurons in the crow pallium whose activity tracks the bird's report of a stimulus rather than the stimulus itself, a neural correlate of sensory consciousness of the kind used in human studies (M7). The common ancestor of birds and mammals was an amniote, and the relevant structures derive from a shared pallial organization. By Definition~\ref{d:verdict}, birds are conscious.
\end{proof}

The homology claimed in the proof is of pallium to pallium, not of layered cortex to avian nuclei; the inference does not require that consciousness depend on cortical lamination, which no human evidence shows. If, as \citet{Merker2007} argues, subcortical structures are sufficient for basic consciousness, the homology is even closer, since those structures are conserved across amniotes.

\section{Cephalopods}

\begin{lemma}\label{l:ceph}
Cephalopods are probably conscious.
\end{lemma}

\begin{proof}
Octopuses, squid and cuttlefish show the marker families with high convergence. Octopuses guard and tend wounds (M3); injected with a noxious substance, they avoid a chamber in which they experienced it and prefer one in which they received a local anesthetic, which also abolishes wound-directed behavior (M4, M5) \citep{Crook2021}; they make flexible trade-offs (M2); they show complex learning, exploration and play \citep{Mather2008} (M5). The review of \citet{Birch2021} found strong evidence of sentience across most criteria. But cephalopod nervous systems evolved independently of the vertebrate one, from a common ancestor with bilaterians that had at most a simple nerve net; the integrative structures are analogous, not homologous (Definition~\ref{d:hom}). By Principle~\ref{p:conv} the inference is strong; by Principle~\ref{p:hom} it does not receive the support of homology. By Definition~\ref{d:verdict}, cephalopods are probably conscious.
\end{proof}

\paragraph{Distributed processing.} The octopus nervous system is remarkable not only for its independent origin but for its organization. Of roughly half a billion neurons, about two thirds are in the arms, and each arm can perform complex movements, such as reaching and fetching, under local control \citep{Hochner2012,GodfreySmith2016}. The central brain issues commands that the arms elaborate and receives information that they have filtered. Does such distributed processing count against consciousness?

\begin{proposition}\label{p:dist}
Distributed processing does not count against the consciousness of cephalopods. It bears on the structure of their consciousness rather than on its presence.
\end{proposition}

\begin{proof}
The markers that support Lemma~\ref{l:ceph} are shown by the animal as a whole: the octopus, not its arm, avoids the chamber and guards the wound, and the anesthetic's effects are global. Integration of information sufficient to produce these effects is therefore present, whatever its anatomical location (M7). The human mind is itself distributed over semi-autonomous processes, as MND-21 holds, and that does not count against human consciousness; what consciousness requires is integration at the level of the subject, not a single central processor. Distributed processing may mean that octopus experience is less unified than ours, or that the arms carry some processing that contributes to the animal's experience only through its results. MND-15 holds that even split-brain patients have partially unified consciousness rather than two streams; on that model, the octopus may have a single, partially unified field whose arm-related contents are less integrated with one another than human bodily contents are.
\end{proof}

The question whether each arm might be a separate subject is not settled by the evidence, and Proposition~\ref{p:dist} does not exclude it. But the evidence that supports the ascription of consciousness is evidence about the animal, so the default is that the animal is the subject. The companion manuscript MND-23-b treats the case in more detail.

\section{Insects}

\begin{lemma}\label{l:insects}
For insects, the evidence remains inconclusive.
\end{lemma}

\begin{proof}
Some markers are present. Bumblebees given an unexpected reward show a dopamine-dependent positive bias in ambiguous situations (M6) \citep{Solvi2016}; fruit flies show nociceptive sensitization and some motivational trade-offs (M1, M2); bees show complex learning, including learning by observation (M5) \citep{Chittka2022}; and the insect central complex performs an integrative function comparable to that of the vertebrate midbrain (M7) \citep{Barron2016}. But other markers are weak, absent or untested. Wound-directed behavior is rarely reported, and early observations of insects continuing to feed while severely injured suggested its absence \citep{Eisemann1984}; evidence of analgesic effects is limited (M3, M4). \citet{Gibbons2022} found the evidence strong in some orders and lacking or unexamined in others. Several positive findings admit explanation by mechanisms of reinforcement learning and global state modulation that need not involve experience. The structures are analogous, not homologous: as for cephalopods, the last common ancestor of insects and vertebrates lies at the split between protostomes and deuterostomes, and had no integrative brain of the relevant kind. By Principle~\ref{p:conv} the inference is weak, and an explanation without consciousness has not been excluded. By Definition~\ref{d:verdict}, the evidence is inconclusive.
\end{proof}

Lemma~\ref{l:insects} does not say that insects are not conscious. The evidence is changing quickly, and the absence of some markers reflects in part the absence of experiments rather than negative results. It is also consistent with Lemma~\ref{l:insects} that consciousness differs across insect orders. The companion manuscript MND-23-c treats the case.

\section{The main result}

\begin{theorem}\label{t:main}
Mammals and birds are conscious; cephalopods are probably conscious, with distributed processing; the evidence for insects remains inconclusive.
\end{theorem}

\begin{proof}
From Lemmas~\ref{l:mammals}, \ref{l:birds}, \ref{l:ceph} and \ref{l:insects}, and Proposition~\ref{p:dist}.
\end{proof}

\begin{remark}\label{r:scope}
Theorem~\ref{t:main} concerns four groups. It does not issue verdicts for fishes, reptiles, amphibians, crustaceans or other invertebrates, for which the method of this manuscript would have to be applied in turn. For vertebrates other than mammals and birds, the homology of subcortical structures gives the inference of Principle~\ref{p:hom} some support; we do not assess the markers here.
\end{remark}

\section{Robustness across theories}\label{s:robust}

The method of this manuscript is theory-light, but its verdicts can be checked against the main theories of consciousness. Let a theory be \emph{admissible} if it is consistent with the human contrast cases from which the markers of Definition~\ref{d:markers} are drawn: it counts as conscious the processes that, in humans, are reportable, flexible and globally available, and as unconscious those, such as blindsight and masked priming, that are not. Four families of admissible theory are now defended. On \emph{global workspace} theories, consciousness is the broadcast of information to many consuming systems \citep{Dehaene2014}. On \emph{recurrent} and integration theories, it is sustained reverberant processing in sensory systems. On \emph{midbrain} theories, it is the integrated model of the animal's body and environment that structures such as the superior colliculus and its homologues construct for the control of action \citep{Merker2007,Barron2016}. Higher-order theories form a fourth family, but MND-14 argues that they fail, since they cannot handle misrepresentation without positing consciousness of non-existent states; we set them aside.

\begin{proposition}[Robustness]\label{p:robust}
For mammals and birds, every admissible theory in the three remaining families yields the verdict of consciousness. For cephalopods, the theories yield it if they are stated functionally, and fall silent if they are stated in terms of vertebrate anatomy. For insects, the theories diverge.
\end{proposition}

\begin{proof}
Mammals and birds have structures homologous with those each family identifies in humans: a pallial system with widespread reciprocal connections, recurrent sensory processing, and midbrain integrative structures (Lemmas~\ref{l:mammals} and \ref{l:birds}). Each family's criterion is satisfied by homologous structures doing homologous work. Cephalopods have no homologues of these structures, but the central brain's vertical and frontal lobes sustain learning and integrate sensory information across modalities, and the markers of Lemma~\ref{l:ceph} indicate global availability of nociceptive information to flexible control. A functional statement of any of the three families is therefore satisfied, while an anatomical statement has nothing to apply to. For insects, the midbrain family is satisfied by the central complex on the argument of \citet{Barron2016}; the workspace family requires a degree of global broadcast whose presence in brains of about a million neurons is in dispute; and recurrent theories require sustained processing whose presence in insect sensory systems is not established. The theories give different answers.
\end{proof}

Proposition~\ref{p:robust} explains the pattern of Theorem~\ref{t:main} from another direction. Where the verdict is that a group is conscious, the verdict does not depend on which admissible theory is true. Where it is that a group is probably conscious, the residual doubt is whether consciousness depends on anatomy or on function, a doubt that the analogy-based inference of Principle~\ref{p:hom} records. Where the evidence is inconclusive, the theories disagree, and the markers do not yet discriminate between them. The verdicts are thus robust to theoretical uncertainty in exactly the cases in which they are firm.

\section{Further cases}

\paragraph{Unremarkable species.} The verdicts concern groups, and the evidence most often cited, from primates, corvids and octopuses, concerns their most intelligent members. Definition~\ref{d:cons} is not a claim about intelligence, and the markers that matter most for it, M1 to M4, are shown by species with no reputation for cognitive sophistication. Lame broiler chickens select feed containing an analgesic, and their gait improves when they do \citep{Danbury2000}; laboratory rodents show the full set of markers. Since the structures concerned are homologous across each class, the verdicts for mammals and birds extend to their typical members, and in particular to the farmed species on which ETH-10 relies.

\paragraph{Self-control in cephalopods.} Cuttlefish will forgo an immediately available, less preferred food in order to obtain a preferred one after a delay, with a tolerance for delay comparable to that of some large-brained vertebrates \citep{Schnell2021}. The finding is evidence of flexible, integrated control (M2, M7) and supports Lemma~\ref{l:ceph}; it adds nothing to the homology side of the inference, and so does not move the verdict from probable to settled.

\paragraph{Play in insects.} Bumblebees roll small balls repeatedly without reward, and do so more when young, which has been interpreted as play and so as a sign of positive affect \citep{Dona2022}. If the interpretation is right, the finding adds to M6. But play-like behavior can be produced by intrinsically rewarding motor routines whose operation need not be experienced, and the finding does not discriminate between the theories that disagree about insects (Proposition~\ref{p:robust}). It is the kind of evidence that, accumulated, could move the verdict, and does not yet do so.

\paragraph{Anesthesia and sleep.} Within each group, the markers lapse under general anesthesia and in deep sleep, and reappear on recovery. This is a further check on the method: the markers are sensitive to the manipulations that abolish consciousness in humans. MND-25 holds that dreams are conscious experiences during sleep; that some mammals and birds show sleep stages resembling human REM sleep is consistent with the verdict for them, though it is not used in its support.

\section{Recent literature}

The marker approach has been criticized from two sides. From the restrictive side, \citet{Carruthers2019} argues that if the global workspace theory is correct, consciousness is a matter of global broadcast, and the question whether an animal's global broadcast is sufficiently like ours to count as conscious has no determinate answer. We reply that Principle~\ref{p:ibe} does not depend on any one theory: the markers are those that all major theories associate with consciousness in humans, and the inference is to the best explanation of their joint presence. If no theory-neutral question about animal consciousness had an answer, neither would the ethical questions of ETH-07, and we see no reason to accept that consequence. From the expansive side, \citet{Mikhalevich2020} argue that the restriction of moral concern to vertebrates rests on bias, and \citet{Tye2017} argues that the best explanation of behavioral similarity is consciousness for a wide range of animals, including some insects and crustaceans. Theorem~\ref{t:main} agrees for cephalopods and holds that for insects the explanation by consciousness has not yet been shown to be better than its rivals.

\citet{Birch2020} distinguish dimensions along which consciousness may vary between species: perceptual richness, evaluative richness, integration at a time and across time, and selfhood. Proposition~\ref{p:dist} is a claim about one of these dimensions, integration at a time, and illustrates their point that the question \enquote{which animals are conscious?} is followed by the question \enquote{conscious in what ways?}. \citet{Birch2024} proposes that, where the evidence for sentience is a realistic possibility, precautionary measures proportionate to the risk are warranted. That proposal concerns action under uncertainty and is compatible with Lemma~\ref{l:insects}: an inconclusive verdict is one under which precaution may be appropriate.

\section{Objections and replies}

\begin{objection}[Morgan's canon]\label{o:morgan}
All the markers can be explained by unconscious mechanisms. Morgan's canon requires that we prefer such explanations, so consciousness should be ascribed to no animal.
\end{objection}

\begin{reply}
The canon, applied consistently, would deny consciousness to other human beings, whose behavior can also be described as the output of mechanisms; MND-06 explains why that conclusion is unwarranted. For animals closely related to us, cladistic parsimony favors shared consciousness (Principle~\ref{p:hom}), and the canon must be weighed against it \citep{Sober2005}. For distant groups, the canon has force, which is why the verdicts of Theorem~\ref{t:main} weaken with evolutionary distance.
\end{reply}

\begin{objection}[No cortex]\label{o:cortex}
Consciousness in humans depends on the cortex. Animals without a layered cortex, including birds, cannot be conscious \citep{Key2016}.
\end{objection}

\begin{reply}
The objection assumes that what matters is the anatomical form of the cortex rather than the functions it performs. Birds perform the relevant functions with a pallium homologous with the cortex, and their neurons track report rather than stimulus \citep{Nieder2020}. Children born without most of the cortex show signs of awareness and affect \citep{Merker2007}. The dependence of human consciousness on cortex is a dependence on certain kinds of integration, which other structures can provide.
\end{reply}

\begin{objection}[Theory first]\label{o:theory}
The question which animals are conscious cannot be answered without a theory of consciousness, and the marker approach smuggles in a theory.
\end{objection}

\begin{reply}
The markers are features that, in humans, distinguish conscious from unconscious processing; they are identified by contrast cases within the human case, and they are shared by the main theories as consequences. The approach requires a theory only in the weak sense of a hypothesis about what consciousness does, and it is designed to be robust across the theories now available. A theory that turned out to be correct could revise the verdicts, and Theorem~\ref{t:main} is held subject to that revision.
\end{reply}

\begin{objection}[Bias in favor of the familiar]\label{o:bias}
The markers were chosen from mammalian, and especially human, behavior. They will tend to find consciousness in animals that resemble us and miss it in those that do not. The inconclusive verdict for insects reflects this bias.
\end{objection}

\begin{reply}
The markers are drawn from the human case because it is the only case in which we know independently which processes are conscious. Their mammalian origin does bias the method toward false negatives for very different animals, which is why Lemma~\ref{l:insects} is a verdict of inconclusiveness and not a negative verdict. The cephalopod case shows that the method can find consciousness in animals very unlike us.
\end{reply}

\section{Consequences}

For ETH-07, which holds that sentience is necessary and sufficient for moral status, Theorem~\ref{t:main} implies that mammals and birds have moral status, that cephalopods probably do, and that for insects the question is open. ETH-10 relies on the first conjunct in its argument that factory farming, whose victims are chiefly mammals and birds, is impermissible. ETH-07-b treats the evidence on invertebrate sentience and how to act on it under uncertainty; the verdicts of Theorem~\ref{t:main} are its starting point. For the treatment of cephalopods in research, the probable verdict supports the inclusion of cephalopods in the scope of animal welfare protections, as has been done in several jurisdictions.

\section{Conclusion}

The question which animals are conscious can be answered without a settled theory of consciousness, by inference to the best explanation of convergent markers, weighted by homology. Mammals and birds are conscious: the markers converge and the structures are homologous with ours. Cephalopods are probably conscious: the markers converge, but in an independently evolved and distributed nervous system. For insects, the evidence remains inconclusive.

\begin{thebibliography}{99}\small
\bibitem[Barron and Klein(2016)]{Barron2016} Barron, A.~B. and Klein, C. (2016). What insects can tell us about the origins of consciousness. \emph{Proceedings of the National Academy of Sciences}, 113(18), 4900--4908.
\bibitem[Birch(2022)]{Birch2022} Birch, J. (2022). The search for invertebrate consciousness. \emph{No\^us}, 56(1), 133--153.
\bibitem[Birch(2024)]{Birch2024} Birch, J. (2024). \emph{The Edge of Sentience: Risk and Precaution in Humans, Other Animals, and AI}. Oxford University Press.
\bibitem[Birch et~al.(2021)]{Birch2021} Birch, J., Burn, C., Schnell, A., Browning, H. and Crump, A. (2021). \emph{Review of the Evidence of Sentience in Cephalopod Molluscs and Decapod Crustaceans}. LSE Consulting, London School of Economics and Political Science.
\bibitem[Birch et~al.(2020)]{Birch2020} Birch, J., Schnell, A.~K. and Clayton, N.~S. (2020). Dimensions of animal consciousness. \emph{Trends in Cognitive Sciences}, 24(10), 789--801.
\bibitem[Carruthers(2019)]{Carruthers2019} Carruthers, P. (2019). \emph{Human and Animal Minds: The Consciousness Questions Laid to Rest}. Oxford University Press.
\bibitem[Chittka(2022)]{Chittka2022} Chittka, L. (2022). \emph{The Mind of a Bee}. Princeton University Press.
\bibitem[Crook(2021)]{Crook2021} Crook, R.~J. (2021). Behavioral and neurophysiological evidence suggests affective pain experience in octopus. \emph{iScience}, 24(3), 102229.
\bibitem[Danbury et~al.(2000)]{Danbury2000} Danbury, T.~C., Weeks, C.~A., Chambers, J.~P., Waterman-Pearson, A.~E. and Kestin, S.~C. (2000). Self-selection of the analgesic drug carprofen by lame broiler chickens. \emph{Veterinary Record}, 146(11), 307--311.
\bibitem[Dehaene(2014)]{Dehaene2014} Dehaene, S. (2014). \emph{Consciousness and the Brain: Deciphering How the Brain Codes Our Thoughts}. Viking.
\bibitem[Dona et~al.(2022)]{Dona2022} Dona, H.~S.~G., Solvi, C., Kowalewska, A., M\"akel\"a, K., MaBouDi, H. and Chittka, L. (2022). Do bumble bees play? \emph{Animal Behaviour}, 194, 239--251.
\bibitem[Edelman et~al.(2005)]{Edelman2005} Edelman, D.~B., Baars, B.~J. and Seth, A.~K. (2005). Identifying hallmarks of consciousness in non-mammalian species. \emph{Consciousness and Cognition}, 14(1), 169--187.
\bibitem[Eisemann et~al.(1984)]{Eisemann1984} Eisemann, C.~H. et~al. (1984). Do insects feel pain? A biological view. \emph{Experientia}, 40, 164--167.
\bibitem[Emery and Clayton(2004)]{Emery2004} Emery, N.~J. and Clayton, N.~S. (2004). The mentality of crows: Convergent evolution of intelligence in corvids and apes. \emph{Science}, 306(5703), 1903--1907.
\bibitem[Feinberg and Mallatt(2016)]{Feinberg2016} Feinberg, T.~E. and Mallatt, J.~M. (2016). \emph{The Ancient Origins of Consciousness: How the Brain Created Experience}. MIT Press.
\bibitem[Gibbons et~al.(2022)]{Gibbons2022} Gibbons, M., Crump, A., Barrett, M., Sarlak, S., Birch, J. and Chittka, L. (2022). Can insects feel pain? A review of the neural and behavioural evidence. \emph{Advances in Insect Physiology}, 63, 155--229.
\bibitem[Godfrey-Smith(2016)]{GodfreySmith2016} Godfrey-Smith, P. (2016). \emph{Other Minds: The Octopus, the Sea, and the Deep Origins of Consciousness}. Farrar, Straus and Giroux.
\bibitem[Hochner(2012)]{Hochner2012} Hochner, B. (2012). An embodied view of octopus neurobiology. \emph{Current Biology}, 22(20), R887--R892.
\bibitem[Jarvis et~al.(2005)]{Jarvis2005} Jarvis, E.~D. et~al. (2005). Avian brains and a new understanding of vertebrate brain evolution. \emph{Nature Reviews Neuroscience}, 6(2), 151--159.
\bibitem[Key(2016)]{Key2016} Key, B. (2016). Why fish do not feel pain. \emph{Animal Sentience}, 1(3).
\bibitem[Mather(2008)]{Mather2008} Mather, J.~A. (2008). Cephalopod consciousness: Behavioural evidence. \emph{Consciousness and Cognition}, 17(1), 37--48.
\bibitem[Merker(2007)]{Merker2007} Merker, B. (2007). Consciousness without a cerebral cortex: A challenge for neuroscience and medicine. \emph{Behavioral and Brain Sciences}, 30(1), 63--81.
\bibitem[Mikhalevich and Powell(2020)]{Mikhalevich2020} Mikhalevich, I. and Powell, R. (2020). Minds without spines: Evolutionarily inclusive animal ethics. \emph{Animal Sentience}, 5(29).
\bibitem[Morgan(1894)]{Morgan1894} Morgan, C.~L. (1894). \emph{An Introduction to Comparative Psychology}. Walter Scott.
\bibitem[Nagel(1974)]{Nagel1974} Nagel, T. (1974). What is it like to be a bat? \emph{The Philosophical Review}, 83(4), 435--450.
\bibitem[Nieder et~al.(2020)]{Nieder2020} Nieder, A., Wagener, L. and Rinnert, P. (2020). A neural correlate of sensory consciousness in a corvid bird. \emph{Science}, 369(6511), 1626--1629.
\bibitem[Schnell et~al.(2021)]{Schnell2021} Schnell, A.~K., Boeckle, M., Rivera, M., Clayton, N.~S. and Hanlon, R.~T. (2021). Cuttlefish exert self-control in a delay of gratification task. \emph{Proceedings of the Royal Society B}, 288, 20203161.
\bibitem[Sober(2005)]{Sober2005} Sober, E. (2005). Comparative psychology meets evolutionary biology: Morgan's canon and cladistic parsimony. In L. Daston and G. Mitman (eds.), \emph{Thinking with Animals: New Perspectives on Anthropomorphism}. Columbia University Press.
\bibitem[Solvi et~al.(2016)]{Solvi2016} Solvi, C., Baciadonna, L. and Chittka, L. (2016). Unexpected rewards induce dopamine-dependent positive emotion-like state changes in bumblebees. \emph{Science}, 353(6307), 1529--1531.
\bibitem[Tye(2017)]{Tye2017} Tye, M. (2017). \emph{Tense Bees and Shell-Shocked Crabs: Are Animals Conscious?} Oxford University Press.
\end{thebibliography}

\end{document}
